\documentclass{homework}
% \usepackage{graphicx} % Required for inserting images
\usepackage{amsmath}
\usepackage{gensymb}
\usepackage{subfig}
\usepackage{float}

\class{ECEN 5478}
\title{Homework 4}
\author{Sean Jennings}
\date{October 12, 2023}

\begin{document}
\maketitle

\section{Problem 1}

\renewcommand{\N}{\mathcal{N}}
\begin{letterenum}
\item Let $G_t := \mathcal{I} - \nabla f_t$ and $\hat{G}_t := G_t + \N_t$
where $\N_t: \R^n \to \R^n$ is the constant map which returns $\nu_t\in\R^n$.
Then, the algorithmic map can be written
\begin{equation*}
    x_t = \hat{G}(x_{t-1})
\end{equation*}
and $\nabla f_t(x_t^*) = 0$ implies that
\begin{equation*}
    x_t^* = G_t(x_t^*)
\end{equation*}

Since $f_t$ is $\mu$-strongly convex and $L$-smooth, from lectures and the previous
homework, we have that $G_t$ is Lipschitz w/ constant
$\rho = \max\{|1 - \alpha L|, |1 - \alpha \mu|\}$. Then, a single-step bound on
tracking error can be derived as follows:
\begin{align}
    \norm{x_t - x_t^*} &= \norm{\hat{G}_t(x_{t-1}) - G_t(x_t^*)} \\
        &= \norm{G_t(x_{t-1}) - G_t(x_t^*)} \\
        &= \norm{G_t(x_{t-1}) + \nu_t - G_t(x_t^*)} \\
        &\leq \norm{G_t(x_{t-1} - G_t(x_t^*))} - \norm{\nu_t} \\
        &\leq \norm{G_t(x_{t-1}) - G_t(x_{t-1}^* + G_t(x_{t-1}^*))}
            + \norm{G_t(x_{t-1}^* - G_t(x_t^*))} + \norm{\nu_t} \\
        &\leq \rho \norm{x_{t-1} - x_{t-1}^*} + \rho \sigma_{t-1} + \norm{\nu_t}
        \label{eq:single-step-bound}
\end{align}
where $\sigma_t := \norm{x_{t+1}^* - x_t^*}$.

Iterating over equation \ref{eq:single-step-bound} yields the $t$-step tracking
bound:
\begin{align}
    \norm{x_t - x_t^*} &\leq \rho^2 \norm{x_{t-2} - x_{t-2}^*} + \rho^2 \sigma_{t-2} +
        \rho \sigma_{t-1} + \rho \norm{\nu_{t-1}} + \norm{\nu_t}
        &\leq \rho^t \norm{x_0 - x_0^*} + D_t + E_t
        \label{eq:multi-step-bound}
\end{align}
where
\begin{equation*}
    D_t = \sum_{i=1}^t \rho^{t - i + 1} \sigma_{i-1}
\end{equation*}
and 
\begin{equation*}
    E_t = \sum_{i=1}^{t} \rho^{t-i} \norm{\nu_i}
\end{equation*}

\item We compute the $\limsup$ of the summands of equation \ref{eq:multi-step-bound}
separately:

\begin{align*}
    \limsup_{t\to\infty} \rho^t \norm{x_0 - x_0^*} &= 0 \\
    \limsup_{t\to\infty} E_t &=
        \limsup_{t\to\infty} \sum_{i=1}^t \rho_{t-i} \norm{\nu_i} \\
        &\leq \limsup_{t\to\infty} \sum_{j=0}^{t-1} \rho^j \hat{\nu} \\
        &= \frac{\hat{\nu}}{1-\rho} \\
    \limsup_{t\to\infty} D_t &=
        \limsup_{t\to\infty} \sum_{i=1}^{t} \rho^{t-i+1} \sigma_{i+1} \\
        \leq \limsup_{t\to\infty} \biggl(
            \sum_{j=0}^{t-1} \rho^{t+1}
        \biggr) \sup_{\tau \geq 0} \sigma_\tau \\
        &= \rho \hat{\sigma} \limsup_{t\to\infty} \sum_{j=0}^{t} \rho^t \\
        &= \frac{\rho \hat{\sigma}}{1 - \rho}
\end{align*}

where $\hat{\nu} = \sup_{t \geq 0} \norm{\nu_t}$ and
$\hat{\sigma} = \sup_{t \geq 0} \norm{sigma_t} =
\sup_{t \geq 0} \norm{x_{t+1}^* - x_t^*}$. (We note a slight disprepency between
this definition of $\hat{\sigma}$ and the one in the problem statement which uses
$x_{t-1}^*$ rather than $x_{t+1}^*$.)

Then the $\limsup$ of the tracking bound is:
\begin{equation*}
    \limsup_{t\to\infty} \norm{x_t - x_t^*} =
        \frac{\hat{\nu}} + \rho \hat{\sigma}{1 - \rho}
\end{equation*}
\end{letterenum}

\section{Problem 3}

\begin{letterenum}
\item Let $f_t(x) = 1/2 \norm{Ax - w_t}^2$ where $f_t: \R^2 \to \R$,
$w_t \in \R^2$, $A\in \R^{2\times 2}$ and additionally
$w_t = [w_{t,1}, w_{t,2}]^T$ has
\begin{align*}
    w_{t,1} = 5 + \sin(t/10)
    w_{t,2} = 2 + \sin(t/100)
\end{align*}
and
\begin{equation*}
    A = \mat{1 &0.8 \\ 0.7 & 1}
\end{equation*}

Let $\{x_t\}_{t \in \mathbb{N}}$ be the sequence generated by the online gradient
method, i.e.
\begin{equation*}
    x_{t+1} = x_t - \alpha \nabla f_t(x_t)
\end{equation*}
where $\nabla f_t(x) = A^T (Ax - w_t)$, and $x_0$ is an input, which we will
arbitrarily select to be $x_0 = [10, 5]^T$.

To select the step size, we note that $\nabla^2f_t(x) = A^T A$. Since
$\nabla^2 f_t(x)$ is constant and positive definite,
it is $\mu$-strongly convex and $L$-smooth for all $x\in\R^2$ and $t\in\R$.
Furthermore, we can compute convexity and smoothness parameters, $\mu$ and $L$
as follows (as proven in the previous homework set):
\begin{align*}
    \mu = \lambda_{min} (\nabla^2 f(x)) = \lambda_{min} (A^T A) \approx 0.0631 \\
    L = \lambda_{max} (\nabla^2 f(x)) = \lambda_{max} (A^T A) \approx 3.067 \\
\end{align*}

Defining $G_t := \mathcal{I} - \alpha \nabla f$, from the previous homework,
we have that the $G_t$ is Lipschitz with constant
$\rho(\alpha) = \max \{ | 1- \alpha L|, |1 - \alpha \mu| \}$.

We proceed to compute the convergence bound for our algorithm. 
so that $x_{t+1} = G_t(x_t)$, we have
\begin{align*}
    \norm{x_t - x_t^*} &= \norm{G_{t-1}(x_{t-1} - x_t^*)} \\
        &= \norm{G_{t-1}(x_{t-1}) - G_{t-1}(x_{t-1}^*)} +
            \norm{x_{t-1^*} - x_t^*} \\
        &= \rho(\alpha) \norm{x_{t-1} - x_{t-1}^*} + \sigma_{t-1}
\end{align*}
where $\sigma_t = \norm{x_{t+1}^* - x_t^*}$ is the path length at time $t$.

The multistep bound is then:
\begin{align}
    \norm{x_t - x_t^*} &= \rho^t \norm{x_0 - x_0^*} +
        \sum_{i=0}^{t-1} \rho^{t-i-1}\sigma_i
        &= \rho^t \norm{x_0 - x_0^*} + \frac{1 - \rho^t}{1-\rho} \hat{\sigma}
        \label{eq:multistep-bound-alg}
\end{align}
with $hat{\sigma} = \sup_{\tau \geq 0} \sigma_\tau$.

Since the only dependence upon $\alpha$ within the tracking bound is via
$\rho(\alpha)$, we can minimize the tracking bound by selecting $\alpha$ to
minimize $\rho$:
\begin{equation*}
    \rho_{min} = \min_\alpha > 0 \rho(\alpha)
\end{equation*}
Since $\mu < L$, the maximum occurs when $-(1-\alpha L) = 1 - \alpha \mu$:
\begin{align*}
    -(1 - \alpha L) &= 1 - \alpha \mu \\
    \alpha (L + \mu) &= 2  \\
    \alpha = \frac{2}{L + \mu}
\end{align*}

 Since the given
$A$ is invertible and our optimization problem is unconstrained, our
fixed point occurs when $\nabla f_t(x) = A^T(Ax - w_k) = 0$, so that
\begin{align*}
    x_t^* = A^{-1} w_t
\end{align*}

Given our knowledge of $w_t$ we can bound the path length.
\begin{align}
    \sigma_t &= \norm{x_{t+1}^* - x_t^*} 
        &= \norm{A^{-1} w_{t+1} - A^{-1} w_t}
        &\leq \norm{A^{-1}} \norm{w_{t+1} - w_t}
        \label{eq:norm}
\end{align}

The difference between consequentive $w_t$s is given by:
\begin{equation*}
    w_{t+1,1} - w_{t,1} = t + \sin(\frac{t+1}{10}) - (5 + \sin{\frac{t}{10}})
        = \sin(\frac{t}{10} + \frac{1}{10}) - \sin(\frac{t}{10}) 
\end{equation*}
which can be bounded by converting into complex coordinates. For any
$\omega,\phi \in \R$, we have
\begin{align*}
    \sin(\omega t + \phi) - \sin(\omega t) 
        &= \Im (e^{i(\omega t + \phi)}) - \Im (e^{i\omega t}) \\
        &= \Im (e^{i\omega t} e^{i\phi} - e^{i\omega t}) \\
        &= \Im ( e^{i\omega t} (e^{i\phi} - 1)) \\
        &\leq | e^{i\omega t} ( e^{i\phi} - 1)| \\
        &\leq | ( e^{i\phi} - 1)|
        &= \sqrt{(\cos \phi - 1)^2 + \sin^2 \phi}\\
        &= \sqrt{2 ( 1- \cos \phi)}
\end{align*}
So that the magnitude on the right-hand side of equation \ref{eq:norm} is
bounded by:
\begin{equation}
    \hat{w} := \sup_{t \geq 0} \norm{w_{t+1} - w_t} =
        \sqrt{4(1 - \cos(1/10))^2 + 4(1 - \cos(1/100))^2} \\
        = 2 \sqrt{(1 - \cos(1/10)^2 + (1 - \cos(1/100))^2)}
        \label{eq:disturbance-bound}
\end{equation}
and the bound for path length is:
\begin{equation}
    \hat{\sigma} \leq = \norm{A^{-1}} \hat{w}
        \label{eq:path-length-bound}
\end{equation}

\item The algorithm has been coded in Python and plots of the results
are shown in figures \ref{} - \ref{}.

Figure \ref{} shows the tracking error as a function of the iteration
index, with the tracking error shown on both linear and logarithmic
scales. Additionally, we plot the upper bound on tracking error
as given by equations \ref{eq:multistep-bound-alg}, \ref{eq:disturbance-bound},
and \ref{eq:path-length-bound}. We note that in this case the upper bound
is rather conservative.

% \begin{center}
%     \includegraphics[]{}
% \end{center}

Interestingly, when $A$ is selected instead to be the identity matrix,






\end{letterenum}


\end{document}
